\section*{Chapter \ref{ch:b2:michelson} --- The Michelson Interferometer}

\begin{solution}{exo:b2:michelson:1}
$N = 2\Delta e/\lambda = 340$; $250$ fringes: $\Delta e = 125\lambda = \qty{74}{\micro m}$; a tenth of a
fringe is $\lambda/20 = \qty{30}{nm}$.
\end{solution}

\begin{solution}{exo:b2:michelson:2}
$p_0 = 2e/\lambda = 1697.8$: not an integer, the centre is not bright ($\varepsilon =
0.8$). $r_k = f\sqrt{(k - 1 + \varepsilon)\lambda/e}$: $3.1$, $4.6$, \qty{5.7}{cm}. At \qty{0.05}{mm} the
radii grow by $\sqrt{10}$: few, broad rings.
\end{solution}

\begin{solution}{exo:b2:michelson:3}
$\lambda/2\alpha = \qty{3.2}{mm}$; six fringes across \qty{2}{cm}; doubling $\alpha$ halves the
spacing. Contact: translate the mirror until the white-light fringes
(a white one flanked by coloured ones) appear --- they exist only within
a micrometre of $\delta = 0$.
\end{solution}

\begin{solution}{exo:b2:michelson:4}
$2(n - 1)e/\lambda = 177$ fringes; $0.2$ fringe gives $\Delta n = 0.2\lambda/2e = \qty{6e-4}{}$.
A centimetre of glass adds millimetres of dispersive path: the colours
no longer agree anywhere, no white fringe.
\end{solution}

\begin{solution}{exo:b2:michelson:5}
(a) $\lambda^2/2\Delta\lambda = \qty{0.291}{mm}$. (b) $5.00/17 = \qty{0.294}{mm}$: $\Delta\lambda = \lambda^2/2\Delta e =
\qty{0.590}{nm}$; $0.1\%$: $\pm\qty{0.0006}{nm}$. (c) At $e = \qty{0.145}{mm}$ (the first
disappearance): $p_1 = 492.4$, $p_2 = 491.9$ --- half an order apart; at
\qty{0.29}{mm} they differ by one. (d) The weaker fringes survive: $V_{\min} =
(I_1 - I_2)/(I_1 + I_2) = 1/3$.
\end{solution}

\begin{solution}{exo:b2:michelson:6}
(a) $2(n - 1)\ell = 94\lambda$: $n - 1 = \qty{2.77e-4}{}$. (b) $0.94$ fringe per percent;
$\dd(n - 1)/\dd T = -(n - 1)/T = \qty{-9.5e-7}{K^{-1}}$. (c) A few millibars or a
degree in one arm move the fringes: hence vacuum or equal arms. (d)
$0.1/94$: $\pm\qty{3e-7}{}$.
\end{solution}

\begin{solution}{exo:b2:michelson:7}
(a) $\ell_c = \lambda^2/\Delta\lambda = \qty{3}{cm}$: rings fade beyond $e \approx \qty{1.5}{cm}$. (b)
\qty{0.3}{mm}, $e \approx \qty{0.15}{mm}$. (c) $\ell_c = \qty{300}{m}$: yes. (d) $\ell_c \approx \qty{1}{\micro m}$:
$2e < \qty{1}{\micro m}$, a couple of orders.
\end{solution}

\begin{solution}{exo:b2:michelson:8}
(a) $2(n - 1)e_g = \qty{5.2}{mm}$; yes, by \qty{2.6}{mm} of mirror travel. (b) The
glass path varies with $\lambda$ by $2e_g\Delta n = \qty{100}{\micro m} \gg \ell_c$: no
wavelength-independent zero. (c) $2(n - 1)\Delta e_g < \qty{0.3}{\micro m}$: $\Delta e_g <
\qty{0.3}{\micro m}$ --- a relative tilt below $\qty{1.5e-5}{rad}$ over \qty{2}{cm}. (d) A
pellicle is micrometres thick; a cube is symmetric: both beams see
the same glass.
\end{solution}

\begin{solution}{exo:b2:michelson:9}
(a) $\Delta t \approx Lv^2/c^3 = \qty{4.1e-15}{s}$, a path of \qty{1.2}{\micro m}. (b) Swapping the
arms doubles it: $2Lv^2/c^2\lambda = 0.4$ fringe. (c) $v < 30\sqrt{0.01/0.4} =
\qty{4.7}{km/s}$. (d) No motion through an ether is seen; the speed of light
is the same along both arms --- the postulate of relativity.
\end{solution}

\begin{solution}{exo:b2:michelson:10}
(a) $I = S_0[1 + \cos(4\pi\nu_0e/c)]$, period $\lambda/2$ in $e$. (b) Each frequency
gives a fringe of slightly different period; a Gaussian band of width
$\Delta\nu$ sums to fringes with a Gaussian envelope of width $\sim c/2\Delta\nu$.
(c) $R = \nu/\Delta\nu = 2e_{\max}/\lambda = 2 \times 10^5$. (d) All wavelengths fall on
one detector at once (Fellgett) through a wide aperture (Jacquinot):
far more signal than a slit spectrometer in the photon-starved
infrared.
\end{solution}

\begin{solution}{exo:b2:michelson:11}
(a) Each arm's path changes by $2N\Delta L$, in opposite senses: $\Delta\delta = 4N\Delta L
= \qty{4.8e-15}{m}$, $\Delta\varphi = 2\pi\Delta\delta/\lambda = \qty{2.8e-8}{rad}$. (b) Near a dark
fringe the detector sees an intensity proportional to $\Delta\varphi^2$ --- or,
with a small offset, to $\Delta\varphi$: a relative change of order $10^{-8}$. (c)
$5 \times 10^{20}$ photons per second, $5 \times 10^{17}$ per millisecond, relative
noise $1.4 \times 10^{-9}$: the signal is twenty times the noise; more power
lowers the noise as $1/\sqrt P$. (d) $\Delta L \propto L$; the pendulum suspensions
filter the ground's motion above their resonance.
\end{solution}

\begin{solution}{exo:b2:michelson:12}
(a) For a point $P$ at depth $e(P)$ and incidence $i$, $\delta = 2e(P)\cos i$
plus a term $\propto\alpha i$ from the tilted second reflection, which
vanishes at normal incidence: only there do both rays from $P$ reach
the eye --- localization on the mirrors. (b) $ei^2 < \lambda/4$: $i_{\max} <
\sqrt{\lambda/4e} = \qty{0.09}{rad}$. (c) The rings depend on $i$ alone: any eye
position sees the same angular pattern. (d) The wedge viewed off the
mirrors; focus on $M_2$ or reduce the tilt.
\end{solution}

\begin{solution}{pb:b2:michelson:1}
\textbf{1.} Wedge fringes; $\lambda/2\alpha = \qty{2}{mm}$: $\alpha = \qty{1.6e-4}{rad}$.

\textbf{2.} $M_1'$ has become parallel to $M_2$: the air plate, whose
\omterm{prop:b2:two-wave-interference:film}{fringes of equal inclination} are rings.

\textbf{3.} $158\lambda/2 = \qty{50}{\micro m}$; rings swallowed at the centre: $e$
decreasing.

\textbf{4.} Contact, $e = 0$. Sodium light has no feature there (one
wavelength, uniform field either way); white light shows the one
position where every colour has $\delta = 0$, flanked by colours as the
orders separate.

\textbf{5.} $\delta = 2\alpha x$: a white fringe on the contact line and one or
two coloured ones on each side ($2\alpha x < \ell_c \approx \qty{1}{\micro m}$).

\textbf{6.} $r_{10} = f\sqrt{10\lambda/e} = \qty{5}{mm}$: $e = 10\lambda f^2/r^2 = \qty{1.0}{cm}$.

\textbf{7.} Parallel, the two beams cross equal glass at every
wavelength; without $C$, the dispersion of \qty{5}{mm} of glass leaves no
common zero: no white fringe.

\textbf{8.} $\lambda^2/4\Delta\lambda = \qty{0.145}{mm}$; then every $\lambda^2/2\Delta\lambda = \qty{0.291}{mm}$.

\textbf{9.} $11$ intervals over \qty{3.20}{mm}: \qty{0.291}{mm}, $\Delta\lambda = \qty{0.597}{nm}$;
$\pm\qty{10}{\micro m}$ on \qty{3.2}{mm}: $\pm\qty{0.002}{nm}$.

\textbf{10.} $492.4$ and $491.9$ (half an order apart); at the return,
$984.7$ and $983.7$.

\textbf{11.} $1/3$.

\textbf{12.} $\ell_c = \lambda^2/\Delta\lambda = \qty{1.7}{cm}$: $2e < \qty{1.7}{cm}$, about $29$.

\textbf{13.} Three fringe systems: full contrast every $\lambda^2/2\Delta\lambda$, with
two partial dips between --- the pattern of three slits.

\textbf{14.} $2(n - 1)\ell = \qty{29}{\micro m}$: $46$ fringes.

\textbf{15.} $n - 1 = 46\lambda/2\ell = \qty{2.91e-4}{}$.

\textbf{16.} $9.2$ fringes.

\textbf{17.} Equal arms: the drift is common to both and cancels; a
\qty{1}{cm} inequality gives $2 \times 0.01 \times 3 \times 10^{-6} = \qty{0.1}{fringe}$ ---
harmless for both parts.

\textbf{18.} They are present, unchanged, throughout the count: a
constant offset.

\textbf{19.} $71$ fringes; the count gives $n - 1$, a signature of the gas.

\textbf{20.} $31\,600$ fringes over a centimetre, read to a tenth:
\qty{30}{nm}, $3 \times 10^{-6}$; the screw gives a micrometre.

\textbf{21.} $\delta$ swings by $\pm\qty{100}{nm}$, $\pm0.16$ fringe at \qty{100}{Hz}: the
eye sees lowered contrast; at $25$ frames per second the camera
samples four periods per frame and sees a frozen or slowly beating
pattern.

\textbf{22.} $10^{-8} \times 31\,600 = 3 \times 10^{-4}$ fringe: nothing.

\textbf{23.} $\lambda/20 = \qty{32}{nm}$ on \qty{10}{mm}: $3 \times 10^{-6}$.

\textbf{24.} It compares a length directly with a wavelength, fringe by
fringe; the grating spreads all wavelengths at once over a wide range.

\textbf{25.} Length: rings counted at the centre. Doublet: the
periodic loss of contrast. Index: fringes passing while the cell
fills. \omterm{prop:b2:scalar-light-model:coherence}{Coherence length}: the thickness at which the rings fade.
\end{solution}
